\documentclass[dvipdfmx, uplatex]{jsarticle}
\usepackage[dvipdfmx]{graphicx}
\usepackage{amsmath}
\usepackage{amssymb}
\usepackage{amsfonts}
\usepackage{amsthm}
\usepackage{bm}
\usepackage{enumitem}

\usepackage{silence}
\WarningFilter*{hyperref}{Token not allowed in a PDF string}
\usepackage[dvipdfmx]{hyperref}
\usepackage{pxjahyper}
\allowdisplaybreaks[4]

\usepackage[left=25truemm,top=20truemm,bottom=20truemm,right=25truemm]{geometry}

\theoremstyle{definition}
\newtheorem*{ex}{例題}

\title{数理コンテスト 事前アナウンスとウォームアップ例題\\
\large ―― 代数・幾何・解析・確率の 4 分野 ――}
\author{}
\date{}

\begin{document}
\maketitle
\vspace{-4zh}

%======================================================================
\section*{コンテストについて}
\addcontentsline{toc}{section}{コンテストについて}
%======================================================================
本コンテストは大学院修了程度の数理素養をお持ちの社会人を対象とした出題です。形式は次のとおり。

\begin{itemize}[leftmargin=1.6em]
\item \textbf{制限時間 2 時間。}本番は\textbf{代数・幾何・解析・確率}の 4 分野から各 1 題，計 4 題を出します。
各題は独立で，どれから解いても構いません。
\item \textbf{全問を解き切る分量ではありません。}得意分野を選び，深く解き進めてください。各題は多数の
小問に細かく分かれており，序盤は易しく後半ほど重くなります。\textbf{部分点が積み上がる}設計なので，
途中まででも得点になります。
\item \textbf{解答形式。}多くの小問は「\,$\sim$ を求めよ／表式で表せ／値を答えよ／オーダーを答えよ\,」の
形です。\emph{正しく計算できていれば表式が一致する}ことで採点します（証明の体裁は問いません）。
\item \textbf{道具立て。}必要な公式・定義は本番の問題文内で与えます。特殊な前提知識は要求しません
（微分積分・線形代数・初等確率の運用力が中心）。
\end{itemize}

\medskip
\noindent\textbf{このプリントの目的。}本番の中身（核心）には触れず，\emph{各分野の入口の手触り}だけを示す
ウォームアップ例題を分野ごとに 2--3 問用意しました。所要は各問 数分程度。\textbf{解答は末尾にまとめて}
あります。まず自力で手を動かし，肩慣らしと得意分野の見極めにお使いください。

\bigskip
%======================================================================
\section{代数（ウォームアップ）}
%======================================================================
固有値・行列式・その摂動など，行列（線形写像）の言葉での計算です。以下，
$A=\begin{pmatrix}2&1\\ 1&2\end{pmatrix}$ とする。

\begin{ex}[代数 1]
$A$ の固有値と固有ベクトルを求めよ。
\end{ex}

\begin{ex}[代数 2]
$B=\begin{pmatrix}0&1\\ 1&0\end{pmatrix}$ とし，微小パラメータ $\lambda$ を入れた $A+\lambda B$ を考える。
行列式 $\det(A+\lambda B)$ を $\lambda$ の\textbf{1 次まで}求めよ（$\det(A+\lambda B)=\det A\,
\bigl(1+\lambda\,\mathrm{tr}(A^{-1}B)+\cdots\bigr)$ を使ってよい）。
\end{ex}

\begin{ex}[代数 3]
単位行列を $I$ として $\det(A+\lambda I)$ を考える。これが固有値 $\mu_i$ を用いて
$\prod_i(\mu_i+\lambda)$ に等しいことを使い，$\det(A+\lambda I)$ を $\lambda$ の多項式として書き下せ。
\end{ex}

%======================================================================
\section{幾何（ウォームアップ）}
%======================================================================
平面ベクトル場の「渦の数」・閉曲線まわりの循環・写像の巻き数など，位相的・幾何的な量です。

\begin{ex}[幾何 1]
平面の各点 $(x,y)$ に矢印 $\bm v=(P,Q)$ を与える。原点を反時計回りに 1 周するとき矢印 $(P,Q)$ の向きが
回る\emph{正味の回転数}（渦の数，反時計回りを正）を，次の 3 つについて求めよ：
\[
    \text{(a)}\ (x,\,y),\qquad \text{(b)}\ (x,\,-y),\qquad \text{(c)}\ (x^2-y^2,\,2xy).
\]
\end{ex}

\begin{ex}[幾何 2]
原点を中心とする単位円 $C$ に沿った線積分（循環）
\[
    \oint_C \frac{-y\,dx + x\,dy}{x^2+y^2}
\]
を求めよ。さらに，これが原点を 1 周だけ囲む\emph{任意の}単純閉曲線で同じ値になることを述べよ。
\end{ex}

\begin{ex}[幾何 3]
整数 $k$ に対し，閉曲線 $\varphi\mapsto(\cos k\varphi,\ \sin k\varphi)$（$\varphi:0\to2\pi$）は原点のまわりを
何回まわるか（巻き数）。$k=2$ と $k=3$ で答えよ。
\end{ex}

%======================================================================
\section{解析（ウォームアップ）}
%======================================================================
微分方程式に冪級数を代入して係数の漸化式を作り，その解がどんな級数になるかを調べます。
微分方程式
\[
    x^2\,y'(x) + y(x) = x
\]
の，$x=0$ で $0$ となる形式的べき級数解 $y(x)=\displaystyle\sum_{n\ge1} a_n x^n$ を考える。

\begin{ex}[解析 1]
上の級数を方程式に代入し，両辺の $x^n$ の係数を比べて，$a_n$ の満たす漸化式を導け
（$a_1$ の値と，$n\ge2$ での $a_n$ と $a_{n-1}$ の関係）。
\end{ex}

\begin{ex}[解析 2]
その漸化式を解いて $a_n$ を $n$ の閉じた式で表せ。
\end{ex}

\begin{ex}[解析 3]
得られたべき級数 $\sum_{n\ge1}a_n x^n$ の収束半径を求めよ（理由も一言）。この級数は $x\neq0$ で
収束するか。
\end{ex}

%======================================================================
\section{確率（ウォームアップ）}
%======================================================================
期待値の定義に沿って計算します。離散確率変数 $X$（取りうる値 $x$，確率 $P(x)$）と任意関数 $g$ に対し
\[
    \langle g(X)\rangle := \sum_{x} g(x)\,P(x)
\]
を期待値と定める（$\sum_x P(x)=1$）。

\begin{ex}[確率 1]
1 ステップの変位 $\Delta$ が値 $+a,\,-a$ を各確率 $\tfrac12$ でとる。定義に従って
$\langle\Delta\rangle$ と $\langle\Delta^2\rangle$ を $a$ で求めよ。
\end{ex}

\begin{ex}[確率 2]
同じ $\Delta$ について，特性関数 $\langle e^{ik\Delta}\rangle=\sum_\delta e^{ik\delta}P(\delta)$ を求めよ。
さらに，独立な $n$ 個の変位の和 $X_n=\sum_{j=1}^n\Delta_j$ について
$\langle e^{ikX_n}\rangle=\prod_{j=1}^n\langle e^{ik\Delta_j}\rangle$（独立性）を使って $\langle e^{ikX_n}\rangle$ を求めよ。
\end{ex}

\begin{ex}[確率 3]
同じ和 $X_n=\sum_{j=1}^n\Delta_j$（各 $\Delta_j$ は独立同分布）について，
$\langle X_n^2\rangle=\sum_{j,l}\langle\Delta_j\Delta_l\rangle$ を出発点に，独立性
（$j\neq l$ で $\langle\Delta_j\Delta_l\rangle=\langle\Delta_j\rangle\langle\Delta_l\rangle$）を用いて
$\langle X_n^2\rangle$ を $n,a$ で求めよ。
\end{ex}

\newpage
%======================================================================
\section*{解答}
\addcontentsline{toc}{section}{解答}
%======================================================================

\noindent\textbf{代数 1.}\quad 固有方程式 $(2-\mu)^2-1=0$ より $2-\mu=\pm1$，すなわち
$\boxed{\mu=1,\,3}$。固有ベクトルは $\mu=3$ で $\boxed{(1,1)}$，$\mu=1$ で $\boxed{(1,-1)}$。

\smallskip
\noindent\textbf{代数 2.}\quad $A+\lambda B=\begin{pmatrix}2&1+\lambda\\ 1+\lambda&2\end{pmatrix}$ ゆえ
$\det(A+\lambda B)=4-(1+\lambda)^2=3-2\lambda-\lambda^2$。1 次までで
\[
    \boxed{\,\det(A+\lambda B)=3-2\lambda+O(\lambda^2)\,}.
\]
（一般式でも $\det A=3$，$A^{-1}=\tfrac13\begin{pmatrix}2&-1\\ -1&2\end{pmatrix}$ より
$\mathrm{tr}(A^{-1}B)=-\tfrac23$，$\det A\cdot\mathrm{tr}(A^{-1}B)=-2$ で一致。）

\smallskip
\noindent\textbf{代数 3.}\quad 固有値 $1,3$ より $\det(A+\lambda I)=(1+\lambda)(3+\lambda)$，すなわち
\[
    \boxed{\,\det(A+\lambda I)=\lambda^2+4\lambda+3\,}.
\]
（代数 2 で $B=I$ とした場合に当たり，$\det A=3$，1 次係数 $\mathrm{tr}(\mathrm{adj}\,A)=4$，$\lambda^2$ 係数 $\det I=1$。）

\bigskip
\noindent\textbf{幾何 1.}\quad 円上で $(x,y)=(\cos\varphi,\sin\varphi)$ とすると，
(a) $(\cos\varphi,\sin\varphi)$ は向きが $\varphi$ ゆえ $\boxed{+1}$；
(b) $(\cos\varphi,-\sin\varphi)$ は向きが $-\varphi$ ゆえ $\boxed{-1}$（鞍点型）；
(c) $(\cos^2\varphi-\sin^2\varphi,\,2\sin\varphi\cos\varphi)=(\cos2\varphi,\sin2\varphi)$ は向きが $2\varphi$ ゆえ $\boxed{+2}$。

\smallskip
\noindent\textbf{幾何 2.}\quad 単位円上 $x=\cos\varphi,\,y=\sin\varphi$ では
$-y\,dx+x\,dy=(\sin^2\varphi+\cos^2\varphi)\,d\varphi=d\varphi$，$x^2+y^2=1$ ゆえ
\[
    \oint_C\frac{-y\,dx+x\,dy}{x^2+y^2}=\int_0^{2\pi}d\varphi=\boxed{2\pi}.
\]
被積分形式は原点以外で「回転なし」（完全形式）なので，原点を 1 周囲む任意の単純閉曲線で値は同じ $2\pi$
（原点を囲まなければ $0$）。これは曲線の\emph{巻き数} $\times2\pi$ である。

\smallskip
\noindent\textbf{幾何 3.}\quad 偏角が $k\varphi$ ゆえ $\varphi:0\to2\pi$ で偏角は $2\pi k$ 増える。巻き数は $k$：
$k=2$ で $\boxed{2}$ 回，$k=3$ で $\boxed{3}$ 回。

\bigskip
\noindent\textbf{解析 1.}\quad $y=\sum_{n\ge1}a_nx^n$，$x^2y'=\sum_{n\ge1}n a_n x^{n+1}=\sum_{m\ge2}(m-1)a_{m-1}x^m$。
$x^2y'+y=x$ の各次を比べて，$x^1$：$\boxed{a_1=1}$；$x^m\ (m\ge2)$：$(m-1)a_{m-1}+a_m=0$，すなわち
\[
    \boxed{\,a_m=-(m-1)\,a_{m-1}\quad(m\ge2)\,}.
\]

\smallskip
\noindent\textbf{解析 2.}\quad $a_1=1$ から $a_m=-(m-1)a_{m-1}$ を反復して
\[
    \boxed{\,a_n=(-1)^{n-1}(n-1)!\,}\qquad(a_1=1,\ a_2=-1,\ a_3=2,\ a_4=-6,\dots).
\]

\smallskip
\noindent\textbf{解析 3.}\quad 比 $\bigl|a_{n+1}/a_n\bigr|=n\to\infty$ ゆえ収束半径 $\boxed{R=0}$。
$x\neq0$ ではどこでも発散する（係数が階乗で増える形式的発散級数）。

\bigskip
\noindent\textbf{確率 1.}\quad 定義より $\langle\Delta\rangle=\tfrac12(+a)+\tfrac12(-a)=\boxed{0}$，
$\langle\Delta^2\rangle=\tfrac12a^2+\tfrac12a^2=\boxed{a^2}$。

\smallskip
\noindent\textbf{確率 2.}\quad $\langle e^{ik\Delta}\rangle=\tfrac12 e^{ika}+\tfrac12 e^{-ika}=\boxed{\cos ka}$。
独立性より $\langle e^{ikX_n}\rangle=\prod_{j=1}^n\langle e^{ik\Delta_j}\rangle=\boxed{\cos^n(ka)}$。

\smallskip
\noindent\textbf{確率 3.}\quad $j\neq l$ では独立かつ平均 $0$ ゆえ $\langle\Delta_j\Delta_l\rangle=0$，$j=l$ では
$\langle\Delta_j^2\rangle=a^2$。よって
\[
    \langle X_n^2\rangle=\sum_{j,l}\langle\Delta_j\Delta_l\rangle=\sum_{j=1}^n a^2=\boxed{n\,a^2}.
\]

\end{document}
